% Pursuing Stacks (cotes 134-2 to 134-8): archive cote 134-4, site page 75 (PDF page 76), typescript page 163.
% typescript. Draft compared: ps3.tex:3828-3893.
% AI-transcribed by Claude Opus 5.5 (claude-opus-5-5, adaptive thinking, effort high), Message Batches API, batch ps2 (msgbatch_01NVwr1Yuo8qJf5rykeGBXPK), 2026-11-12; not yet checked by eye.
\scanpage{134-4/75}
\mspage{163}

Coming back to a general M,W, we have \deleted{associated to thisxpair a} \inserted{defined earlier a}
canonical homotopy structure $h_W$ on M, which we may call ``associated to
W'' - this\xout{wxxx} is also the weak homotopy interval structure on M, gene-%
rated by intervals in M which are W-aspheric over e . This makes sense
at least, provided in M finite products exist. If moreover M satisfies
the conditions b) to d) of page \hand{123} concerning sums and connected compo-%
nents of objects of M, we \over{have}{\unsure{had}} define\hand{d}\hand{(}independently of W\hand{)}the \uline{canonical}
homotopy structure $h_M$ in M , which \over{may}{\unsure{may}} be viewed as the weak homotopy
interval structure generated by all 0-connected intervals in M. It still
seems that in all cases we are going to be interested in, we have the
equality
\[ h_W = h_M \quad . \]
When speaking of homotopy notions in M (such as f and g being homotopic
maps, written
\[ f \sim g \quad , \]
or a map being a homotopism, or an object being contractible) it will be
understood \xout{thatxthis} (unless otherwise stated) that this refers to the
homotopy structure $h_W = h_M$ . In case we should not care to impose \hand{(}\inserted[typed]{otherwise} un-%
needed assumptions which will imply $h_W = h_M$ , we'll be careful when
referring to homotopy notions, to say which structure we are working
\uline{with.}
\margin{(horizontal line drawn by hand across the left margin)}

\hand{6}\over{4}{4}. We recall that a set of arrows $W \subset \mathrm{F\unsure{l}}(M)$ is called \uline{saturated} if it
satisfies the conditions:

a) Identities belong to W .

\over{\illegible{1}}{b}) If f,g are maps and gf exists, then if two among f,g,gf are in
W, so is the third.

c) If $f:X \to Y$ and $g: Y \to X$ are such that gf and fg are in W, so are
f and g.

\xout{The} On the other hand, \uline{strong saturation} for W means that W is
the set of arrows made invertible by the localization functor
\[ M \longrightarrow M_W = W^{-1}M \quad , \]
\xout{Txextroublexwith} or equivalently, \hand{that} W can be described as the set of
arrows made invertible by some functor $M \longrightarrow M'$ . The trouble with strong
saturation is that it is a condition which \over{\illegible{1}}{of}ten is not easy to check
in concrete situations. This is so for instance for the notion of weak
equivalence in (Cat), and the numerous variants defined in terms of
cohomology. Therefore we surely wont impose the strong saturation condi-%

\margin{75}
